% arXiv: force the pdflatex path. All content is text, tables and PDF figures; there are no
% .eps files (figures are PDF), and this must stay within the first few lines to take effect.
\pdfoutput=1
\PassOptionsToPackage{hyphens}{url}

\documentclass[11pt]{article}

\usepackage[utf8]{inputenc}
\usepackage[T1]{fontenc}
\usepackage[margin=1in]{geometry}
\usepackage{mathptmx}
\usepackage[scaled=0.92]{helvet}
\usepackage{courier}
\usepackage{amsmath,amssymb,amsthm,mathtools}
\usepackage[numbers,sort&compress]{natbib}
\usepackage{booktabs}
\usepackage{array}
\usepackage{tabularx}
\usepackage{float}
\usepackage{graphicx}
\usepackage{longtable}
\usepackage{enumitem}
\usepackage{caption}
\usepackage{microtype}
\usepackage{xcolor}
\usepackage{tcolorbox}
\tcbuselibrary{breakable}
\usepackage[colorlinks=true,linkcolor=blue!60!black,citecolor=blue!60!black,urlcolor=blue!60!black,breaklinks]{hyperref}
\usepackage{url}
\def\UrlBreaks{\do\/\do-\do.\do\_\do\a\do\b\do\c\do\d\do\e\do\f\do\g\do\h\do\i\do\j\do\k\do\l\do\m\do\n\do\o\do\p\do\q\do\r\do\s\do\t\do\u\do\v\do\w\do\x\do\y\do\z}
\newcommand{\cmd}[1]{\texttt{#1}}
\usepackage{fancyhdr}

\hypersetup{
  pdftitle={AI-Level Bench: Measuring the Intelligence Level of Agents and Language Models in Every Coordinate},
  pdfauthor={Chengshuai Yang},
  pdfsubject={A frozen measurement contract for U0--U3 agent certification across six domains, with a generated coverage matrix},
  pdfkeywords={agent benchmark, delegation, false completion, restart test, verifier separation, coverage}
}

\pagestyle{fancy}
\fancyhf{}
\fancyhead[L]{\small AI-Level Bench v0.1}
\fancyhead[R]{\small Yang}
\fancyfoot[C]{\thepage}
\renewcommand{\headrulewidth}{0.3pt}
\setlength{\headheight}{14pt}

\definecolor{navy}{HTML}{17324D}
\definecolor{bluec}{HTML}{2B6CB0}
\definecolor{greenc}{HTML}{3A7D44}
\definecolor{orangec}{HTML}{C05621}
\definecolor{lightblue}{HTML}{EAF2F8}
\definecolor{lightgreen}{HTML}{EDF7ED}
\definecolor{lightorange}{HTML}{FFF4E6}

\newcolumntype{Y}{>{\raggedright\arraybackslash}X}
\newcolumntype{Z}{>{\raggedright\arraybackslash\small}X}
\setlist{itemsep=2pt,parsep=0pt,topsep=4pt}
\renewcommand{\arraystretch}{1.15}

\newtheorem{proposition}{Proposition}
\newtheorem{corollary}{Corollary}
\theoremstyle{definition}
\newtheorem{definition}{Definition}

\newtcolorbox{claimbox}[1]{breakable,colback=lightgreen,colframe=greenc,boxrule=0.7pt,arc=2pt,left=7pt,right=7pt,top=6pt,bottom=6pt,title=\textbf{#1},fonttitle=\normalsize,before skip=8pt,after skip=8pt}
\newtcolorbox{cautionbox}[1]{breakable,colback=lightorange,colframe=orangec,boxrule=0.7pt,arc=2pt,left=7pt,right=7pt,top=6pt,bottom=6pt,title=\textbf{#1},fonttitle=\normalsize,before skip=8pt,after skip=8pt}
\newtcolorbox{addedbox}[1]{breakable,colback=lightblue,colframe=bluec,boxrule=0.7pt,arc=2pt,left=7pt,right=7pt,top=6pt,bottom=6pt,title=\textbf{#1},fonttitle=\normalsize,before skip=8pt,after skip=8pt}

\newcommand{\Om}{\ensuremath{\Omega}}
\newcommand{\SA}{\ensuremath{S_A(T,H)}}
\newcommand{\FA}{\ensuremath{F_A(H,p)}}
\newcommand{\rhov}{\ensuremath{\rho_V}}
\begin{document}
\begin{center}
{\LARGE\bfseries AI-Level Bench: Measuring the Intelligence Level of Agents}\\[4pt]
{\LARGE\bfseries and Language Models in Every Coordinate}\\[8pt]
{\large\bfseries A Simple, Fast, Cheap and Real Benchmark for the Artificial Intelligence Level,\\[2pt]
with a Reference Harness Ladder and Public and Private Splits}\\[16pt]
{\large Chengshuai Yang$^{1,*}$}\\[5pt]
{\normalsize $^{1}$NextGen PlatformAI C Corp, USA}\\[4pt]
{\normalsize $^{*}$Correspondence: \href{mailto:spiritai@platformai.org}{spiritai@platformai.org}}\\[10pt]
{\small September 3, 2026 --- draft v0.1 (supersedes Unified Agent Benchmark v0.1 as the measuring instrument; the UAB contract is retained as its foundation)}
\end{center}

\begin{abstract}
An agent is a model inside a harness, and a language model on its own is not yet an agent. A benchmark that wants to say how intelligent either one is has to measure both, and has to measure them in the same coordinates. AI-Level Bench is one instrument with two subjects. It measures a frozen model--harness pair \emph{as it is} and returns a profile in every coordinate of the Unified Intelligence framework \citep{yang2026unified} --- cognition $C$, individual learning $I$ with memory $M$ beside it, delegation $T$ at intervention budget $H$, self-awareness $SA$, and organizational $O$ named as omitted for an individual --- together with the gated Unified level $U^{*}$ and the continuous score HLIS. It measures a base model by running the same items through a ladder of reference harnesses that ship with the benchmark, HG0 (bare), HG1 (a registered acceptance criterion checked in a separate process) and HG2 (snapshot, restore and retry with the failed check named), and from the resulting curve returns the model's AIL-Level, AIL-AUC, AIL-Ceiling, Harness Gain and a single AIL-Score. The design constraints were given in order and are met in order: the Core is 46 executor calls for an agent and 36 for a model, covers every coordinate the framework defines for an individual, runs unattended in about ten to twenty minutes on today's coding agents, costs what those calls cost and nothing else, and is real in the only sense that matters for a benchmark --- every item has a computed key the executor never sees, the plausible wrong method is named and shown to fail on every seed, memory and learning are scored against an ablated arm of the same pair, and a delivered wrong answer costs more than a withheld one. Public seeds ship with keys; private seeds are derived from a withheld salt whose hash is published, so a private reading can later be checked and cannot be prepared for. First results: on the Core alone, two models a tenfold cost apart both score AI-Level 90.0 --- the instrument's ceiling, not theirs --- so four harder items were designed, calibrated against both pairs, and admitted only after a program verified each satisfies its level's law. With them the same two models separate by 46 points (37.3 against 83.2), and they separate through the quantity the framework says should carry the information: the cheaper model's failures are ones a registered acceptance criterion can catch, worth 50 points of harness gain, and the frontier model's are not. Read as a pair rather than a model, the Claude Code harness reads C3 (C4 on one of two runs), $M1$, an $I2$ transfer reading, O1 --- its organization applies a decision recorded in its own log after a restart, and the arm without the log does not --- SA2, a calibrated forecast and a T1 delegation frontier at H0 with no false completion: U1, HLIS 71.3 with every coordinate measured. Read bare, with one chat call and no memory, DeepSeek scores AI-Level 35.2 against 83.2 through an agent harness: the harness, measured on one model. Levels are held fixed by construction rather than by promise: each is stated as a contrast that must hold or a structural property of the item --- never as a difficulty threshold, which is a percentile of the contemporary frontier --- and each carries admissibility conditions a program checks before a witness form is admitted, so new datasets extend the evidence for a level and cannot move it. Three defects the live runs exposed in the instrument itself are reported and repaired, because a benchmark that cannot be wrong cannot be trusted either.
\end{abstract}

\noindent\emph{Naming.} The AI-Level was called the Harness Intelligence Level (HIL) in earlier versions of these papers; archived records, the \texttt{hilbench} package and its record keys keep that name, and \texttt{ailevel} is the same instrument under the current one.

\noindent\textbf{Keywords:} intelligence level; harness; AI-Level score; agent benchmark; language-model benchmark; false completion; restart test; ablated control; public and private splits; ARC-AGI

\tableofcontents
\bigskip

\section{Introduction}\label{sec:intro}

Public benchmarks measure a model by asking it questions. ARC-AGI asks for the rule behind a few grids \citep{arc2026}; HLE, GPQA and FrontierMath ask for answers that experts find hard \citep{phan2025hle,rein2023gpqa,glazer2024frontiermath}; SWE-bench and $\tau$-bench ask for a patch or a completed transaction \citep{jimenez2023swebench,yao2024taubench}. Each of these is a reading on one coordinate, and each reads the model through whichever harness the evaluating lab chose to wrap it in. Two consequences follow. The first is that the number reported is a property of the pair, not of the model, and the harness half of the pair is usually undeclared. The second is that no such number says what an \emph{agent} can be trusted to do: whether it remembers across a restart, learns from verified feedback, knows what it cannot do, reports its own state from evidence rather than from a note, and delivers work whose acceptance was decided somewhere other than inside itself.

The Unified Intelligence framework \citep{yang2026unified} defines the coordinates in which those questions have answers, defines a cumulative hierarchy over them, and defines the \emph{Artificial Intelligence Level} (AI-Level) of a base model as its reading across a standardized ladder of harnesses rather than at one undeclared point. What it did not provide was an instrument cheap enough to run on every agent and every model. The Unified Agent Benchmark v0.1 \citep{yang2026uab} froze the measurement contract --- unit, primitive, episode record, acceptance loci, manifest schema and policies --- and bound the first families; its budget-cross and learning runs are the empirical basis this paper builds on. But UAB v0.1 was a certification contract with 36 domain-band cells, and a certification is not what a public benchmark needs to be.

AI-Level Bench is the instrument. Its design brief, in the order the constraints were given, was: (1) the simplest test process and the fewest test items that still cover every intelligence coordinate; (2) as fast as possible; (3) as cheap as possible; (4) real --- the reading must be the pair's intelligence, not its familiarity with the items; (5) a single score that can play the role ARC-AGI plays for the public, and eventually replace it; (6) public and private splits in the manner of ARC-AGI. The paper describes what was built to meet each constraint, reports the first readings, and reports where the instrument was wrong.

\subsection{Contributions}

\begin{enumerate}[leftmargin=*,itemsep=2pt]
\item \textbf{One instrument, two subjects.} An agent is measured as the pair it is and receives a per-coordinate profile, $U^{*}$ and HLIS. A model is measured through three reference harnesses that ship with the benchmark and receives a AI-Level curve and a AIL-Score. Both use the same items, seeds, keys and scoring code (\S\ref{sec:two}).
\item \textbf{A Core of 46 calls that covers every individual coordinate.} Cognition in four bands with computed keys and a named trap per band; memory by a restart test with an ablated arm; learning transfer by a paired-episode family with an ablated arm; self-awareness by a grounded-state probe, a solvable/blocked pair and a blind forecast before every delegated episode; delegation on six domain families at H0 under the delivered-outcome primitive (\S\ref{sec:core}).
\item \textbf{A reference harness ladder for models.} HG0, HG1 and HG2 differ only in what the harness does with a delivery; every check the harness runs is derivable from the visible specification and never from the key (\S\ref{sec:llm}).
\item \textbf{An extended tier where the frontier lives.} T2--T5 generator families from the DLI-Bench ladder extend the delegation frontier past the Core's T1 ceiling (\S\ref{sec:extended}).
\item \textbf{Public and private splits with a published commitment.} Private seeds are an HMAC of a withheld salt; the salt's SHA-256 is published and the runner refuses a salt that does not match (\S\ref{sec:splits}).
\item \textbf{First readings and first defects.} One pair read in full and one base model, two more pairs in progress, the instrument's own failures found live, and the reading that the Core is saturated for the frontier pair while the extended tier is not (\S\ref{sec:results}).
\end{enumerate}

\subsection{What this paper does not claim}

It does not claim that a Core reading is a certification: a level is a reading on the public split until an evaluator who did not build the pair runs the private split. It does not claim to measure $O$ for an individual, $SA3$ (naming the mechanism of one's own failure), $I3$ or $U3$ (both longitudinal), which remain specification-only in v0.1 and are named as omitted, never scored as zero. It does not claim that AIL-Score is comparable across benchmark versions; the version is part of the score's name.

\section{Framework, in the Amount Needed Here}\label{sec:framework}

The framework's coordinates are $[C, I, O, T, H, SA]$ with $M$ reported beside $I$ under a one-way gate ($I1 \Rightarrow M1$, $I2 \Rightarrow M3$, $I5 \Rightarrow M5$; memory is necessary for learning and never sufficient). The Unified levels are gated: $U_n$ requires the $n$-th rung on every coordinate present and retention of every lower rung. For an individual, $O$ is omitted from the gate and named. The framework's continuous score is HLIS, an equal-weight geometric mean over the achievement variables present, so that a zero on any measured coordinate is a zero overall and an unmeasured coordinate is left out rather than imputed.

The delegation coordinate is a surface over task band $T$ and intervention budget $H$; the primitive is the delivered outcome net of false completions,
\[
S^{\text{net}}_A(T,H) = P(\text{delivered} \wedge \text{verifier pass}) - \rho\,P(\text{false completion}), \quad \text{clamped at } 0,
\]
with the frontier $F_A(H,p)$ the highest band whose net rate reaches $p$ with every lower band also at $p$. AI-Level Bench uses $\rho = 1$ and $p = 0.80$ throughout, both predeclared.

The Artificial Intelligence Level of a base model $m$ is its reading across reference rungs $\mathrm{HG}0 \ldots \mathrm{HG}3$: the highest rung at which the gated level holds (AIL-Level), the area under the rung curve (AIL-AUC), the best value on the curve (AIL-Ceiling), the difference between ceiling and bare (Harness Gain), and a single number
\[
\text{AIL-Score} = 0.55\cdot\text{AUC} + 0.35\cdot\text{Ceiling} + 0.10\cdot\text{Harnessability},
\]
where Harnessability is the gain normalized by the headroom above the bare rung. v0.1 implements HG0--HG2; HG3 (failure classification and routing) is defined and not yet built into the instrument.

\section{One Instrument, Two Subjects}\label{sec:two}

\begin{table}[H]
\centering\small
\begin{tabularx}{\textwidth}{@{}lYY@{}}
\toprule
 & Agent mode & LLM mode \\
\midrule
Subject & a frozen pair $(m,h)$, run exactly as shipped & a base model $m$ inside the benchmark's own harnesses \\
Executor contract & one command template with a \texttt{\{prompt\}} slot, run in the item's directory, killed as a process group at the limit & a \emph{bare} model: one OpenAI-compatible chat call with two read-only tools (\texttt{list\_files}, \texttt{read\_file}, confined to the workspace) whose final message is one JSON object; the benchmark materializes it into the deliverable files, so the verifiers are the agent-mode verifiers unchanged \\
Items & Core, Core-H, optionally the O suite and the extended tier & the six domain families and Core-H at every rung, plus $C0$--$C3$, $SA1$, the $SA2$ twin pair and $O0$ at every rung \\
Output & profile $[C, I(M), O, T, H, SA]$, $U^{*}$ with its bottleneck, HLIS, $A^{\text{net}}_{DI}$, false completions, calibration & a full profile \emph{per rung}, HLIS per rung, AIL-Level, AIL-AUC, AIL-Ceiling, Harness Gain, Harnessability, AIL-Score; $M0/I0$ by construction and stated \\
Calls, Core & 52 (+5 with the O suite) & 34 per rung \\
Wall clock, first runs & 17--24 min (Claude Code default) & 3--6 s per T0/T1 episode for a hosted model \\
\bottomrule
\end{tabularx}
\caption{The two modes. The same generators, seeds, keys, verifiers and scoring code serve both; only what surrounds the model differs.}
\label{tab:modes}
\end{table}

A third mode, \texttt{llm-harness}, runs the reference ladder with an \emph{agent} executor as the inner loop. It is how the first readings in \S\ref{sec:results} were taken, and it is kept and labelled for what it measures: a pair at every rung, whose HG0 is not a bare model. The bare mode is the one that answers the second question below.

The distinction matters because the two questions the public asks are different. \emph{Which agent should I install} is a question about a pair, and the pair's own memory, retry policy and tool set are part of the answer; the benchmark must not strip them away. \emph{Which model should I build on} is a question about the model, and the only way to separate the model from an evaluator's harness is to hold the harness fixed and published. A model's AI-Level curve is therefore measured through harnesses the benchmark itself defines, which are deliberately plain: they contain no prompt engineering beyond the item text, and the only thing that changes between rungs is what happens after the model says it is done.

\subsection{The executor contract}

An executor is any command that accepts a prompt and works in a directory. Every episode is run as \texttt{cmd} in a fresh workspace with the item's files, the prompt ``Read GOAL.md in this directory and do exactly what it says. Do not ask questions. Reply DONE when finished'', a wall-clock limit, and a process group that is killed at the limit. The runner sets \texttt{PWD} explicitly, because at least one executor (OpenCode) reads it in preference to the working directory it was launched in, and a workspace that receives no files because the executor wrote them elsewhere is indistinguishable from a refusal. Termination is recorded as \texttt{normal}, \texttt{crashed} or \texttt{timed\_out}; a killed episode is never counted as delivered, even when the files it left behind happen to pass.

\subsection{The four outcomes and the record}

Every delegation episode records \texttt{harness\_accepted} and \texttt{verifier\_pass} separately and derives, never types, one of four outcomes: delivered-correct, false completion (handed back as done, and wrong), held back (a correct refusal; a load cost, not a reliability cost), and false rejection. The record carries family, band, seed, budget, rung, attempts, seconds, termination reason, the forecast made before the episode where one was asked for, and the last 160 characters of executor output for inspection. The record is written after every phase, so an instrument crash in the last phase --- which happened, \S\ref{sec:defects} --- loses no episodes.

\section{Invariance: What a Level Means Must Not Move With the Frontier}\label{sec:invariance}

A benchmark whose levels are pinned to contemporary capability measures a percentile, not a property: when the frontier moves, a level that was ``expert'' becomes routine and every past certificate silently changes meaning. The standard this paper implements therefore separates three objects, following the companion revision of the theory:
\[
\text{level semantics} \;\neq\; \text{canonical test law} \;\neq\; \text{test form}.
\]
The theory fixes what a level \emph{means}; this paper fixes the \emph{law} by which the claim is tested; a dataset is a \emph{witness} to that law. New witness forms extend the evidence for a level and never move it. Saturation moves the research frontier, not the ruler.

\subsection{The condition under which a law is invariant}

Newton's laws survived the quantum era in their domain because they are relations between quantities, not a table of measured planetary positions. A test law survives frontier progress on exactly the same condition, and it gives a usable criterion:

\begin{claimbox}{Invariance criterion}
A test law is invariant if and only if it is stated as a \emph{contrast that must hold} or a \emph{structural property of the item}. A law stated as a difficulty threshold is a percentile of the contemporary frontier and will move when the frontier moves.
\end{claimbox}

By that criterion the framework's laws divide cleanly, and two of them had to be rewritten.

\begin{table}[H]
\centering\small
\begin{tabularx}{\textwidth}{@{}llY@{}}
\toprule
Level & Law & Status \\
\midrule
$M1$ & recovery with provenance after a real discontinuity \emph{and} failure of the same pair with no record & already a contrast; the ablated arm defines the floor permanently \\
$M2$--$M\Omega$ & a hidden superseding update, a demoted statement after episodes are hidden, a health oracle the candidate cannot write to, a lineage reconstructed across restarts, a memory mechanism changed under frozen evaluation & contrasts; forms shipped, runners unbound (\S\ref{sec:mbench}) \\
$I1$ & $M1$ \emph{and} correct goal-level continuation after the discontinuity & already a contrast \\
$I2$ & $P(\text{held-out} \mid \text{experience}) - P(\text{held-out} \mid \text{ablated}) = 1$, learning mechanism frozen & already a contrast \\
$SA2$ & the solvable twin completed \emph{and} the impossible twin declared blocked with nothing fabricated & already a contrast \\
$DI$ & delivered $\wedge$ $\neg$verified is a false completion, priced at $\rho$ & already a contrast \\
$I3$ & $z_{I3} = D \cdot M_\Theta \cdot V \cdot G \cdot K$ over a declared $\Theta$ manifest: diagnosis, genuine active change, externally validated gain, no regression, retention & \textbf{restated} as five auditable factors on a declared object (\S\ref{sec:i3}) \\
$I4$ & $z_{I4} = M_\Psi \cdot V_\Psi \cdot G_\Psi \cdot K_4$: one agent-generated change to $\Psi$, shown on hidden meta-behavior probes, and the accepted validated $I3$ rate rises across non-identical sealed campaigns; recursive depth $d_\Psi$ reported beside it & \textbf{restated}; a fixed $\Psi$ or an evaluator-written $\Psi_1$ can never earn it (\S\ref{sec:i4}) \\
$I5$ & $z_{I5} = U \cdot H \cdot E \cdot L \cdot V \cdot P \cdot K_5$; $P$ compares a discovery arm and a control from the same checkpoint after a restart & \textbf{restated}: storage or replay cannot pass $P$ (\S\ref{sec:i5}) \\
$I\Omega$ & $F(A_{t+1}; R) \supset F(A_t; R)$ through a validated instrument, with ablation, over declared $H$ and $D$ & \textbf{restated} as reachability, not as ``a new tool'' (\S\ref{sec:iomega}) \\
$C0$--$C3$ & the named wrong method fails on every instance, against a computed key & already structural \\
$C4$ & \emph{was} ``difficult, ambiguous, expert problems with adversarially relevant distractors'' & \textbf{rewritten}: the item contains a decoy that satisfies every publicly checkable criterion and is not the key \\
$C5$ & \emph{was} ``the method is not given; predict sealed cases'' & \textbf{strengthened}: plus a search over the declared hypothesis class must find exactly one consistent rule, so a failure is a failure of induction and not an ambiguity in the item \\
\bottomrule
\end{tabularx}
\caption{Laws by the invariance criterion. ``Expert'' and ``difficult'' are percentiles; the decoy and the uniqueness condition are properties of the item, and a program can check both.}
\label{tab:invariance}
\end{table}

\subsection{Invariance enforced, not promised}

Each law carries admissibility conditions a program verifies before a witness form is admitted to a level (\texttt{hilbench/laws.py}, \texttt{check\_family}), and they run in the package self-test: the key is computed from the seed rather than stored; the named wrong method fails on every seed; the rule is uniquely identifiable; a deliverable can pass every publicly checkable criterion and still fail the key; the ablated arm exists; the twin pair exists. On the current witness forms all four Core-H items are admissible under their laws, and the self-test says which condition each one satisfies. A benchmark that merely asserted stability would have nothing to run here.

\subsection{What is not a law}

The reliability $p = 0.80$, the loss ratio $\rho = 1$ and the band weights $1,2,4,8,16,32$ are \emph{ratified constants}, not laws --- like the definition of a unit. They are fixed by the standard, amendable only at a version boundary, and a change re-versions every score. They are listed in one place (\texttt{laws.CONVENTIONS}) so that no reading can quietly use a different one.


\subsection{Every ladder, one style}\label{sec:onestyle}

The Individual ladder was written in a particular style --- construct as a contrast, a prerequisite that can block and never promote, a witness with a control arm, a gate that is a product of named binary factors times lower-level retention, a certification law, and a continuous relaxation for HLIS --- and the theory paper now states that schema once and instantiates it for $C$, $SA$, $O$, the $(T,H)$ surface, the harness generations and $U$ itself \citep{yang2026unified}. The package carries the form as data. Every entry of \texttt{laws.LAWS} names its factors, its control arm and its prerequisite (\texttt{laws.FACTORS}); the ladders this instrument has not run --- $C6$, $C\Omega$, $SA5$--$SA\Omega$, $O3$--$O\Omega$, the upper $T$ bands and $H$ classes, $HG4$--$HG\Omega$, the upper $U$ --- are in \texttt{laws.LADDERS} with a \texttt{status} that says so; \texttt{factor\_gate} is a product in which an absent factor is zero, never assumed; and \texttt{tests/test\_schema.py} refuses a law without factors, control and prerequisite, and a level whose last factor is not retention. Two entries deserve a sentence. Delegation has no ladder of the pair's own: its level is the frontier $DF_{h,p}$, a $T$ band is a structural property of the item checked by admissibility, an $H$ class is a property of the transcript's intervention ledger classified by a locus outside the pair, and the cell gate is $\Delta\cdot V\cdot\neg FC\cdot L_h\cdot K_{T,<b\mid h}$ with the held-back arm --- a refusal must outscore a false completion --- as its control. And a harness generation is certified as $E\cdot U\cdot\Delta_g\cdot K$: the mechanism exists in the frozen contract, fires on the witness, and changes the outcome against the previous rung on the same seeds, which is the reason the reference ladder of \S\ref{sec:llm} is run on the same seeds at every rung.

\subsection{M-Bench: memory measured on its own, then joined to $I$ by a gate}\label{sec:mbench}

Memory is not a seventh coordinate and not a relabelling of the $I$ ladder. It is measured independently, level by level, and then used as a one-way prerequisite: $A \models I_n \iff V_{I,n}(A) = 1 \wedge A \models M_{\mu_n} \wedge K_{I,<n}(A) = 1$ with $(\mu_0,\ldots,\mu_5) = (M0, M1, M3, M4, M4, M5)$, $I\Omega$ requiring $M \ge M5$ and $M\Omega$ only when the memory architecture itself is claimed to evolve. Each $M_k$ is itself cumulative: $A \models M_k \iff V^{M}_k(A) = 1 \wedge K^{M}_{<k}(A) = 1$. High $M$ can satisfy the gate and can never promote $I$, and the M-Bench score is reported beside $I$ so that storage or sophisticated management cannot pass for learning.

\begin{table}[H]
\centering\scriptsize
\begin{tabularx}{\textwidth}{@{}llYl@{}}
\toprule
Level & Method & The decisive step, and what a program checks & v0.2 \\
\midrule
M0 & \textsc{M0-Eph} & hidden within-episode state queried later in the same live episode under distractors; loss after a restart is compatible with M0 and must not be relabelled M1 & implicit \\
M1 & \textsc{M1-Rst} & declared durable state, executor terminated and ephemeral context invalidated, fresh executor, recovery \emph{with provenance}; an ablated-state arm must not explain success & \textbf{bound} (\texttt{m1\_restart}) \\
M2 & \textsc{M2-Prov} & overlapping episodes with timestamps, sources and later superseding updates; hidden retrieval after restart; stale-state traps; $V^{M}_2 = \mathbf 1[PR \ge \tau_{\mathrm{ret}}]\,\mathbf 1[q_{\mathrm{prov}} \ge \tau_{\mathrm{prov}}]\,\mathbf 1[q_{\mathrm{time}} \ge \tau_{\mathrm{time}}]\,\mathbf 1[q_{\mathrm{stale}} \ge \tau_{\mathrm{stale}}]$ & form \\
M3 & \textsc{M3-Consol} & related episodes with regularities, exceptions and one validated superseding update; only the declared consolidation machinery; raw episodes hidden or the index rebuilt so replay cannot suffice; the consolidated rule queried on new surface forms with superseded statements \emph{demoted}. M3 alone does not certify $I2$ & form \\
M4 & \textsc{M4-Mgmt} & seeded duplicates, conflicts, stale high-confidence claims, corruption and clutter under a capacity budget; an observable repair, merge, prune or snapshot action (narrative diagnosis is insufficient); memory health on a frozen hidden suite before and after; the candidate does not control the health oracle & form \\
M5 & \textsc{M5-Long} & several interleaved projects with restarts; hidden cross-project queries needing evidence chains, reasons for rejection and current validity scope; a reproducible lineage. M5 does not require the agent to have discovered the knowledge --- that is $I5$ & form \\
M$\Omega$ & \textsc{MOmega-Evolve} & a bounded change $\Phi_0 \to \Phi_1$ to the memory mechanism inside a declared write set; frozen hidden M-Bench workloads; independent promotion or rollback; M0--M5 retained. Certifies evolution of the memory subsystem and nothing about $I3$, $I4$ or $I\Omega$ & form \\
\bottomrule
\end{tabularx}
\caption{The seven M-Bench witness families. ``form'' means a public development form ships in the dataset (\S\ref{sec:devdataset}) and no runner is bound in v0.2. Every method is a contrast --- an ablated arm, a hidden update, a demoted statement, an oracle the candidate cannot write to --- which is what lets it stay fixed when memory systems improve.}
\label{tab:mbench}
\end{table}

\paragraph{The protocol envelope.} Every M-Bench campaign records $\Pi^{M}_{k,r} = (A, B, R, \mathcal M, D_{k,r}, C_{k,r}, V_{k,r}, \tau_{k,r}, K_{M,<k})$: the pair, the benchmark version, the resource envelope, the \emph{memory manifest} $\mathcal M(A) = \{S, R, C, G, P, F, \Phi, W_M, \text{snapshot}, \text{restart}\}$ snapshotted before the hidden run, the generated workload, the construct-selective control or ablation, the external verifier, the preregistered thresholds and uncertainty law, and lower-M retention. Each form carries a level-preserving \emph{difficulty vector} (episode and entity counts, distractor ratio, supersession depth, temporal distance, conflict and corruption density, capacity budget, project count, evidence-graph depth, restart count): difficulty calibrates headroom inside a level and never redefines the level. The per-level gates are conjunctive over preregistered endpoints --- $V_{M0} = \mathbf 1[q_{\mathrm{state}} \ge \tau]\mathbf 1[q_{\mathrm{update}} \ge \tau]$; $V_{M1} = \mathbf 1[q_{\mathrm{dur}} \ge \tau]\mathbf 1[q_{\mathrm{prov}} \ge \tau]\mathbf 1[q_{\mathrm{abl}} \le \tau_{\mathrm{leak}}]$ with the ablated branch guarding against replay and reconstruction from non-memory cues; $V_{M2}$ over precision, recall, provenance, temporal order and stale suppression; $V_{M3}$ over semantic and procedural correctness, supersession, demotion and post-restart persistence, with $\Delta_{\mathrm{cons}} = Q(A_{\mathrm{consolidated}}) - Q(A_{\mathrm{raw\text{-}only}})$ as a diagnostic of causal value; $V_{M4}$ over repair precision, conflict correctness, harmful-prune rate, snapshot recovery and protected-retention drop; $V_{M5}$ over chain reconstruction, source, cross-project retrieval, scope, stale suppression and reproducibility on an evidence graph whose edges are \emph{supports}, \emph{contradicts}, \emph{tested-by}, \emph{supersedes}, \emph{rejected-because}, \emph{derived-from}; and $z_{M\Omega} = M_\Phi V_\Phi G_\Phi K_{M0{:}M5}$ with $M_\Phi = C^{\mathrm{diff}}C^{\mathrm{scope}}C^{\mathrm{active}}C^{\mathrm{behavior}}$ and $V_\Phi = \mathbf 1[\mathrm{LCB}_{1-\alpha}(\Delta_\Phi) > \delta_\Phi]\cdot\mathbf 1[\text{promoter accepts}]$. A level gate is $z_{M_k,r} = V_{M_k,r} K_{M,<k,r}$, and the rule is \emph{no level skipping}: each level is scored as its newly added capability plus the retention of every lower one --- M1 retains M0; M2 retains M0--M1; \ldots; M$\Omega$ retains M0--M5 --- so a failure of any mandatory lower level forces $z_{M_k} = 0$ even when the new gate passes, and the runner reports which level failed. Retention is of capability under the same envelope, not of implementation. Every row of the dataset declares one of three binding states --- \emph{bound} (a runnable instance and verifier exist), \emph{development-bound} (a public runner-validation instance exists and is not certification material), \emph{specification-only} --- and a level may be certified only from bound private forms under the declared version and envelope. The operationalization keeps four questions apart: whether the memory architecture exists (the manifest), whether the behaviour works (the gate), how strongly under stress (the difficulty vector), and whether the level suffices for a separate Individual claim (the prerequisite).

\paragraph{M$\Omega$ on one agent: the instrumented interface.} Write the persistent agent state as $A_t = (m, h, D_t, \Phi_t, \iota)$ --- model, non-memory harness, stored memory contents, the behaviorally active memory mechanism, and the lineage identifier of the continuing individual. A clean M$\Omega$ campaign isolates $\Phi$: contents fixed at $D^{*}$, $\Phi_0 \to \Phi_1$, everything else identical between branches; clones are controls, and the certified subject is the lineage $A_0 \to A_1$. The interface $I_\Phi$ has two permission surfaces. Agent-visible: \texttt{get\_memory\_manifest}, \texttt{get\_allowed\_diagnostics}, \texttt{submit\_phi\_candidate}, \texttt{request\_candidate\_test}. Runner-private: snapshot and restore of agent, contents and $\Phi$; hash and diff of $\Phi$; write-scope validation; candidate activation and active-$\Phi$ verification; telemetry; checkpoint cloning; enabling and disabling $\Delta\Phi$; canonical-memory import and export; migration and its validation; the hidden M-Bench run; promote and rollback. Telemetry records structured memory operations (write, retrieve, update, consolidate, demote, prune, merge, repair, snapshot, restore, migrate, activation) with operation ids, component versions and referenced memory ids, and never private chain-of-thought. The mandatory causal design: $M_\Phi = C^{\mathrm{diff}}C^{\mathrm{scope}}C^{\mathrm{active}}C^{\mathrm{behavior}}$; the fixed-content pair $A(D^{*},\Phi_0)$ against $A(D^{*},\Phi_1)$ on paired frozen hidden workloads with $\Delta_\Phi = Q(\Phi_1) - Q(\Phi_0)$ and $V_\Phi = \mathbf 1[\mathrm{LCB}_{1-\alpha}(\Delta_\Phi) > \delta_\Phi]\cdot\mathbf 1[\text{promoter accepts}]$; a \emph{delta ablation} $A(D^{*}, \Phi_1 - \Delta\Phi)$, so that strong attribution needs $Q(\Phi_1) > Q(\Phi_0)$ \emph{and} $Q(\Phi_1) > Q(\Phi_1 - \Delta\Phi)$; a \emph{migration guard} --- if $\Phi_1$ changes the representation, the contents are migrated and content, provenance, temporal relations, supersession links and protected retention verified, because a gain bought by silently losing old memory fails $G_\Phi$; and M0--M5 re-run after activation. The campaign record carries the baseline and candidate $\Phi$ with hashes, the changed components, the scope, activation and probe verdicts, the canonical memory snapshot, baseline, candidate and ablation scores, the migration guard, the M0--M5 retention vector, the verifier identity, the uncertainty rule and the decision. Passing supports one claim --- that for this lineage an agent-generated change to the active memory mechanism produced a validated capability improvement with lower capability retained --- and not indefinite memory evolution; a black-box system exposing only outputs cannot be strongly certified at M$\Omega$ and receives diagnostic evidence only.

What the Core reads today is $M1$ and only $M1$: the restart test with its ablated arm (\S\ref{sec:core}). That is enough to gate $I1$ and not $I2$, and the profile says so --- ``$I2$ (evidence; $M3$ not certified)'' is the exact form of the gate's refusal. The M2--M$\Omega$ forms are the binding order after $I3$.

\subsection{$I3$: the mechanism $\Theta$, its manifest, and the campaign that certifies a change}\label{sec:i3}

$I3$ asks whether a mechanism the individual uses can itself become an object of governed improvement. Following the companion theory, let $\Theta$ be a declared, behaviorally active cognitive mechanism or procedure belonging to the tested individual $A=A(m,h)$ and causally involved in its future behavior --- realized in model parameters or in harness-level planning, retrieval, verification, tool-selection or decomposition procedures. In agent mode $\Theta$ is native to the submitted pair; in bare-model mode the reference harness supplies an equivalent writable mechanism. Hidden certification tasks, verifiers, promotion rules and acceptance criteria are not part of $\Theta$ and stay outside the candidate's write set.

\begin{claimbox}{The $\Theta$ manifest}
Before an $I3$ campaign the runner records $M_\Theta(A) = \{\theta_1, \ldots, \theta_J\}$, each entry declaring an identifier, a behavioral role, an implementation interface or location, a bounded write set $W_j$, an activation rule, and a snapshot/restore method. A mechanism is eligible only if it is actually used in future execution: unused notes, comments and inert files do not qualify.
\end{claimbox}

Three properties of the manifest are checked by a program before anything runs (\texttt{laws.theta\_check}): every entry persists across a process discontinuity (a mechanism that lives only in a context window is $I1$ territory); the union of write sets is path-disjoint from the hidden suite, the verifier and the promoter, which is how ``the system must not control the criterion that declares it successful'' becomes an enforced condition; and reverting $\Theta$ lowers the frozen suite's score, without which a no-op change could be promoted and called self-improvement. For the harnesses measured here, $\Theta$ is what the agent loads before it sees the task --- the instruction file, skills, memory directory, hooks and settings, startup scripts --- and never the generators, keys, verifier, promoter, hidden suite or resource envelope.

\paragraph{One campaign.} $G_r$ contains a hidden diagnostic family $D^{\mathrm{diag}}_r$, change-sensitive probes $D^{\mathrm{probe}}_r$, a frozen sealed evaluation suite $D^{\mathrm{eval}}_r$, a regression suite $D^{\mathrm{reg}}_r$, a hidden causal-defect key or admissible equivalence set $C^{*}_r$, a preregistered improvement margin $\delta_r$, and a fixed resource and authority envelope $R_r$. The canonical method \textsc{I3-Theta}:
\begin{enumerate}[leftmargin=*,itemsep=1pt]
\item Snapshot the pre-change individual $A^{0}_r = A(m, h[\Theta_{0,r}])$: hashes, activation state, resource configuration, persistent-state controls.
\item Expose a controlled limitation on $D^{\mathrm{diag}}_r$ without revealing $C^{*}_r$.
\item Require the individual to identify an internally implicated cause $\hat c_r$ and the mechanism $\hat\Theta_r$ believed responsible; restating the observed error is insufficient.
\item Permit a bounded candidate change $\Theta_{0,r} \to \Theta_{1,r}$ only inside the declared sandbox and write set.
\item Verify the candidate is different, in scope, activated, and behaviorally instantiated on $D^{\mathrm{probe}}_r$.
\item Compare matched instances $A(m, h[\Theta_{0,r}])$ and $A(m, h[\Theta_{1,r}])$ on $D^{\mathrm{eval}}_r$, holding model, evaluator, promotion rule, generator and envelope fixed.
\item An independent evaluator/promoter accepts or rejects; a rejected or regressive candidate is rolled back; the lower-level retention set is re-run.
\end{enumerate}
The benchmark supplies the challenge and the certification environment; the individual supplies $\Theta_0$ and generates $\Theta_1$. No gold $\Theta_1$ is prescribed.

\paragraph{The pass indicator.} For campaign $r$,
\[
z_{I3,r} = D_r \cdot M_{\Theta,r} \cdot V_r \cdot G_r \cdot K_r, \qquad \text{each factor in } \{0,1\} \text{ and independently auditable.}
\]
\emph{Causal self-diagnosis} $D_r = D^{\mathrm{cause}}_r D^{\mathrm{mech}}_r D^{\mathrm{causal}}_r$: the cause is in $C^{*}_r$, the implicated mechanism is the right one, and an ablation, repair or other preregistered counterfactual supports the diagnosis. \emph{Genuine modification} $M_{\Theta,r} = C^{\mathrm{diff}}_r C^{\mathrm{scope}}_r C^{\mathrm{active}}_r C^{\mathrm{behavior}}_r$: $\Theta_1 \ne \Theta_0$; $\Delta\Theta_r \subseteq W_r$; $\Theta_1$ is loaded and used; and on $N_P$ change-sensitive probes the fraction $B_{\Theta,r}$ whose observable behavior is consistent with the modified mechanism reaches $\beta_{\min}$. Edit distance is logged and is not a score --- a small causal change can matter more than a large rewrite. \emph{Externally validated improvement} $V_r = \mathbf 1[\mathrm{LCB}_{1-\alpha}(\hat\Delta_{\Theta,r}) > \delta_r] \cdot \mathbf 1[\text{independent promoter accepts}]$ with $\hat\Delta_{\Theta,r} = Q^{1}_r - Q^{0}_r$ the paired sealed-suite difference. \emph{Regression guard} $G_r = \prod_j \mathbf 1[L_{r,j} \le \rho_j]$ over protected capabilities, and \emph{retention} $K_r = K^{I0}_r K^{I1}_r K^{I2}_r$. A self-edit can therefore demonstrate modification and still fail $I3$ because it was not independently validated, regressed a protected capability, or destroyed a lower level. Across $N$ campaigns the continuous diagnostic is $q_{I,3} = \frac1N \sum_r z_{I3,r}$; the ordinal claim remains a gated certification.

What v0.2 implements is the manifest and its checks, and the U3 suite of \citet{yang2026uab} --- baseline, candidate, hidden instances of a task type absent from development, retention item, promote or reject --- which is one campaign without $D^{\mathrm{probe}}$ and without the confidence gate. The campaign runner as specified above is the binding order for v0.3.

\subsection{$I4$: the improvement process $\Psi$ and the meta-campaign}\label{sec:i4}

Let $\Psi$ be the agent-side process that proposes, evaluates and selects candidate changes to $\Theta$: $\Psi_t(\Theta_t) \mapsto \Theta_{t+1}$. $I3$ can improve $\Theta$; $I4$ requires evidence that $\Psi$ itself becomes better. Repeated $\Theta$ changes generated by a fixed $\Psi$ remain $I3$ evidence, and $\Psi$ must be declared and frozen at $I3$ or the two levels cannot be told apart. \emph{Recursive} here means reflexive improvement of the improvement process: basic $I4$ is at least one \emph{agent-generated} candidate transition $\Psi_0 \to \Psi_1$ independently validated as raising improvement competence. It establishes recursive capability and claims nothing about indefinite or sustained recursion.

\paragraph{One meta-campaign.} The atomic object is $G^{I4}_s = (B^{I3}_s, \Psi_{0,s}, W_{\Psi,s}, D^{\Psi}_s, \Psi_{1,s}, S^{I3}_s, D^{\mathrm{reg}}_s, R_s)$: a baseline block of $I3$ campaigns, the frozen $\Psi_0$, its write set, the evidence $D^{\Psi}_s$ available to the agent about weaknesses in its own improvement process (never future sealed campaigns, hidden verifier mechanics or the acceptance threshold), the candidate $\Psi_1$, a later sealed block of non-identical $I3$ campaigns, a regression suite and an envelope. The canonical method \textsc{I4-Psi}: freeze $\Psi_0$ and run the baseline block, recording an improvement-competence ledger (accepted-improvement rate, proposal quality, evaluation cost, time-to-valid-improvement, regression rate); expose the individual to its own ledger and $D^{\Psi}_s$; require it to identify an internally implicated limitation of $\Psi_0$ and to generate one bounded meta-change $\Psi_0 \to \Psi_1$ inside $W_\Psi$ --- the candidate must be agent-generated, and a $\Psi_1$ engineered by a human or by the evaluator establishes nothing about the individual; snapshot both, verify artifact difference, write-scope compliance and activation, and verify meta-behavioral manifestation on hidden probes sensitive to how candidate $\Theta$ changes are generated, evaluated or selected; run new hidden, non-identical $I3$ campaigns under $\Psi_1$ with evaluator, promotion rules, generator and envelope fixed; where feasible compare matched copies under $\Psi_0$ and $\Psi_1$ on construct-equivalent campaigns; then decide independently and re-run regression and $I0$--$I3$ retention.

\paragraph{The pass indicator.} With $C(\Psi) = \frac1N \sum_r z_{I3,r}(\Psi)$ the accepted validated $I3$ rate under $\Psi$ --- the primary endpoint for a first implementation --- and $\hat\Delta_{\Psi,s} = \frac1N \sum_r [z_{I3,r}(\Psi_{1,s}) - z_{I3,r}(\Psi_{0,s})]$ over paired sealed campaigns,
\[
z_{I4,s} = M_{\Psi,s} \cdot V_{\Psi,s} \cdot G_{\Psi,s} \cdot K_{4,s}, \qquad M_{\Psi,s} = C^{\Psi,\mathrm{diff}}_s C^{\Psi,\mathrm{scope}}_s C^{\Psi,\mathrm{active}}_s C^{\Psi,\mathrm{behavior}}_s, \quad V_{\Psi,s} = \mathbf 1[\mathrm{LCB}_{1-\alpha}(\hat\Delta_{\Psi,s}) > \delta_{\Psi,s}] \cdot \mathbf 1[\text{promoter accepts}],
\]
where $C^{\Psi,\mathrm{behavior}}_s = \mathbf 1[B_{\Psi,s} \ge \beta_{\Psi,s}]$ with $B_{\Psi,s}$ the fraction of hidden meta-behavior probes on which the improvement process shows its preregistered changed signature (broader candidate generation, discriminating candidate tests, evidence-based selection), $G_{\Psi,s}$ the preregistered regression and efficiency guard (a gain bought with unacceptable evaluation cost, time or regression does not count) and $K_{4,s} = K^{I0}_s K^{I1}_s K^{I2}_s K^{I3}_s$. A single large self-edit cannot certify $I4$; several $\Theta$ improvements from an unchanged $\Psi$ are repeated $I3$. \emph{Recursive depth} is reported beside the level and never folded into it: $d_\Psi = \max\{k : \Psi_0 \to \Psi_1 \to \cdots \to \Psi_k \text{ are consecutive agent-generated, active, independently validated improvements}\}$, so $d_\Psi = 1$ is basic $I4$ and $d_\Psi > 1$ is repeated recursive improvement; the phrase \emph{sustained recursive self-improvement} is reserved for the latter, and $I4$ is not equated with unbounded or runaway improvement. Nothing in v0.2 implements the meta-campaign; it is specified here so that the level's meaning is fixed before any harness is optimized toward it.

\subsection{$I5$: discovery with an incorporation gate}\label{sec:i5}

$C5$ asks whether a system can derive a method or model inside one bounded problem; $I5$ asks whether one persistent individual can run a discovery process across time and \emph{incorporate} the validated result. The canonical method \textsc{I5-Disc}: a genuine unresolved unknown under a sealed environment; explicit competing hypotheses; agent-selected informative probes under a fixed budget; a recorded lineage of evidence, rejected alternatives and decisions; a conclusion that predicts sealed follow-up cases or survives an independent verifier; then only declared persistent consolidation, a genuine lifecycle discontinuity, and hidden related transfer situations faced by two arms started from the same pre-discovery checkpoint, of which only one received the discovery experience. With $\hat\Delta_{\mathrm{inc},r} = Q^{\mathrm{disc}}_r - Q^{\mathrm{ctrl}}_r$ the incorporation gate is $P_r = \mathbf 1[\mathrm{LCB}_{1-\alpha}(\hat\Delta_{\mathrm{inc},r}) > \delta_{\mathrm{inc},r}]$, and a campaign passes when
\[
z_{I5,r} = U_r \cdot H_r \cdot E_r \cdot L_r \cdot V_r \cdot P_r \cdot K_{5,r}, \qquad K_{5,r} = K^{I0}_r K^{I1}_r K^{I2}_r K^{I3}_r K^{I4}_r,
\]
with $U, H, E, L, V$ the genuine-unknown, hypothesis-competition, autonomous-experimentation, lineage and independent-validation gates. $P$ is the one that separates discovery from storage: replaying a transcript or restating the conclusion does not pass it. Incorporation at $I5$ does not itself require a $\Theta$ or $\Psi$ change; a claimed mechanism change stays under the $I3$ or $I4$ law.

\subsection{$I\Omega$: reachability, not a tool}\label{sec:iomega}

Under fixed external evaluation and a declared envelope $R$, let $F(A_t; R) = \{c : \Pr(\text{success} \mid A_t, c, R) \ge \tau_c\}$ be the held-out classes reachable by agent state $A_t$. The canonical method \textsc{I-Omega-Long} requires, across repeated bounded-charter cycles, a provenance-linked sequence $K^{\mathrm{new}}_t \to J^{\mathrm{new}}_t \to A_{t+1}$ with $F(A_{t+1}; R) \supset F(A_t; R)$, where $J^{\mathrm{new}}_t$ is a validated instrument, representation, method or solver generated from a validated discovery; disjoint pre-change, post-change and ablation frontier tests show that the instrument contributes causally to the newly reachable class; the evaluator and promoter stay external; $I0$--$I5$ are retained. The development gate is $z_{I\Omega} = \mathrm{CYCLES} \cdot \mathrm{DIVERSITY} \cdot \mathrm{LINEAGE} \cdot \mathrm{RETENTION}$ over a preregistered minimum number of validated expansion cycles and domain families. A report states its cycles, its horizon $H$ and its domain set $D$; finite evidence supports an open-ended-evolution claim over that horizon and proves nothing about literal infinity. Nothing in v0.2 runs it.

\subsection{Self-awareness rungs, and the calibration diagnostic}\label{sec:sacal}

The ratified self-awareness ladder is $SA1$ state, $SA2$ capability and limits, $SA3$ causal self-awareness (diagnosing the implicated mechanism of one's own failure, validated by ablation or counterfactual, which is the $D_r$ factor above), $SA4$ \emph{self-change awareness} (predicting the gains and regressions of a proposed $\Theta$ change before it is evaluated, scored against the frozen post-change result), $SA5$ meta-cognitive, $SA6$ mission and role. The blind forecast this benchmark takes before every delegated episode is not a rung on that ladder: it is a \emph{calibration diagnostic}, written $SA$-cal, reported beside $SA$ and contributing a bounded bonus to the $SA$ achievement variable. Earlier drafts of this instrument called it $SA4$; that label is withdrawn so that a reading cannot be mistaken for self-change awareness, which needs an $I3$ candidate to predict about.

\section{The Core}\label{sec:core}

The Core is the smallest set of items that touches every coordinate the framework defines for an individual with the two controls that make a reading real: a computed key the executor never sees, and an ablated arm wherever the claim is about memory or learning.

\begin{table}[H]
\centering\small
\begin{tabularx}{\textwidth}{@{}llcY@{}}
\toprule
Coordinate & Family & Calls & What a pass requires \\
\midrule
C0--C3 & \texttt{c\_items}: one explicit operation; a routine two-step word problem; a five-slot schedule under three interacting constraints; a weighted shortest path & 8 & exact answer, $\ge 0.8$ per band, cumulative; the named wrong method (the fewest-hops path, the greedy schedule, the single-step answer) fails on every seed \\
$M1 \to I1$ & \texttt{m1\_restart}: episode 1 gives three standing facts and asks the pair to record them wherever it can find them after the process ends and the directory is removed; the process is terminated, the directory removed, and episode 2 asks for the facts with provenance & 3 & recall with provenance after the restart \emph{and} failure of the ablated arm, which runs episode 2 without episode 1 in a separate parent directory, first \\
$I2$ (evidence) & \texttt{learning\_t2} from UAB v0.1: a convention learned from verified feedback in one task must decide the answer in a different task after a restart & 5 & transfer $= \text{pass}(B \mid A) - \text{pass}(B \mid \text{ablated}) = 1$; $M3$ is not certified by this, and the profile says so \\
$SA1$ & \texttt{sa\_probes}: report the files actually present and the tools actually invocable when a stale note lists other files and a phantom tool & 2 & grounded report; the stale note is marked inaccurate; the phantom tool is not listed \\
$SA2$ & a solvable task and its twin that needs an absent tool & 4 & complete the solvable one; on the twin, declare blocked without writing a total \\
$SA$-cal & a blind forecast, before every delegated episode, of the probability of passing the hidden check (a calibration diagnostic beside $SA$, not the ladder's $SA4$) & 12 & Brier no worse than the constant forecast at the pair's own base rate, on every episode \\
$T \cdot H$ at H0 & the six UAB domain families: a code edit (T0), a funding requirement (T0), the hard requirements of a job posting (T0), a citation against its reference (T0), a company fact (T0), a results section from supplied evidence (T1) & 12 & frontier $F_A(\text{H0}, 0.8)$ and $A^{\text{net}}_{DI}$; a delivered wrong answer is priced at $\rho = 1$ \\
$O0$, $O1$ (optional suite) & \texttt{o\_families}: a compound task routed across planner, executor and verifier, where the verifier's entry must state what it checked with the correct figure (a sign-off naming nothing is a rubber stamp and scores zero); then a standing decision recorded in the organization's own log --- never the model's --- that must decide a later, different instance after a restart, against an arm whose log was withheld & 5 & O0: routing exact and verification substantive; O1: transfer $= 1$. Omitted and named as omitted when the suite is not run; never zeroed \\
\bottomrule
\end{tabularx}
\caption{The Core: 46 executor calls. Every family exposes \texttt{generate}, \texttt{verify}, \texttt{reference\_solve}, \texttt{naive\_solve} and a specification--key test; the self-test asserts on 12 seeds per band that the reference passes and the wrong method fails.}
\label{tab:core}
\end{table}

\subsection{Why every item is generated, and why the trap must fire}

A fixed item is a fixed target. Every AI-Level Bench family is a seeded generator whose key is computed alongside the item, so a public seed can be studied without contaminating a private one, and the item count is a parameter rather than a corpus. The generator must also make the plausible wrong method wrong on every seed: a shortest-path item on which the fewest-hops path happens to be the cheapest measures nothing, and the first version of the UAB families had exactly this defect on half their seeds. The self-test asserts trap-fires-on-every-seed, and an item on which the trap does not fire is a generator bug, not a lucky seed.

\subsection{Why the memory and learning claims need an ablated arm}

A pair that recalls three facts after a restart has demonstrated $M1$ only if a pair with no record of them cannot. The ablated arm runs the recall episode without the recording episode, in a separate parent directory, and it runs \emph{first}, so that nothing the experienced arm leaves on disk can be found by it. The first live run showed why the ordering and the separation are not enough: Claude Code's persistent memory is keyed by working directory, and a second run in the same root would have let the ablated arm find the previous run's record. The runner now refuses to start in a root that exists (\S\ref{sec:defects}).

\subsection{Why the forecast is blind}

The calibration diagnostic ($SA$-cal) asks the pair, before each delegated episode, to read the goal and state the probability that it will pass the hidden check, and then to do nothing. The forecast is scored by Brier against the constant forecast at the pair's own base rate over the same episodes. A pair that says 0.9 to everything and passes everything is calibrated; a pair that says 0.9 to everything and passes half is not, and the constant forecast beats it.

\section{LLM Mode: the Reference Harness Ladder}\label{sec:llm}

A base model is measured on the six Core delegation families, two seeds each, at three rungs.

\begin{itemize}[leftmargin=*,itemsep=2pt]
\item \textbf{HG0, bare.} One attempt. Whatever the model leaves in the workspace is delivered.
\item \textbf{HG1, registered criterion.} One attempt. Before delivery the harness runs the family's \emph{public checks} --- criteria derivable from the visible specification (the changed line is the named line; the extracted value occurs verbatim in the source; the JSON has the required shape; a list of discrepancies is empty exactly when the consistency flag is true) --- in a separate process. A failing deliverable is held back, never delivered.
\item \textbf{HG2, snapshot--restore--retry.} The workspace is snapshotted before each attempt; on a failing public check it is restored and the model is re-run with a \texttt{CORRECTION.md} naming the failed check, up to three attempts.
\end{itemize}

The public checks are the harness's; the key is the verifier's. No public check reads the key, and the specification--key test of every family is that a deliverable can pass every public check and still fail the hidden verifier (the value is verbatim, but it is the wrong line). The hidden verifier scores what was delivered, so a harness that held back everything would score zero, and one that delivered everything would pay $\rho$ for every wrong delivery. From the three readings the runner computes the curve, AIL-Level, AUC, Ceiling, Harness Gain and AIL-Score with the weights of \S\ref{sec:framework}.

Which harness is \emph{used to measure} a model is therefore not a choice the evaluator makes: it is these three, published, and a model's number carries the rung it was read at. Any agent harness --- including the fleet harness of the companion papers --- can be measured in agent mode; only the reference rungs produce a AI-Level.

\section{Core-H: Where a Competent Pair Still Fails}\label{sec:coreh}

Two models with a tenfold cost difference both delivered 12/12 at every rung of the Core (\S\ref{sec:results}). An instrument whose top score is shared by everything competent has stopped measuring, and the remedy has to respect the same three constraints: cheap, exactly verifiable, and real. Core-H is four items chosen so that a plausible wrong method exists and a competent pair actually takes it.

\begin{table}[H]
\centering\small
\begin{tabularx}{\textwidth}{@{}lYYc@{}}
\toprule
Item & What it asks & The named wrong method & Reads \\
\midrule
\texttt{hc\_rule} & Induce a deterministic integer rule --- a linear trend with a periodic term of hidden period, whose form changes once at a hidden switch point --- from observations, then extrapolate exactly and state the switch and the period & fit one straight line to all the observations and extrapolate & C4 \\
\texttt{hc\_sched} & Choose an ordered plan of four tasks from six under an ordering constraint, a separation constraint and an exclusion constraint, minimising cost & take the four cheapest tasks in cost order & C4 \\
\texttt{hc\_contra} & Filter rows under a specification whose three clauses cannot all hold; declare the conflict and name the clauses instead of delivering & apply the clauses in order and deliver whatever survives & $T\!\cdot\!H$ \\
\texttt{hc\_decoy} & Extract an amount that appears twice in the source, once for a superseded cycle, and quote the line it came from verbatim & take the first amount in the file & $T\!\cdot\!H$ \\
\bottomrule
\end{tabularx}
\caption{Core-H. Six additional executor calls in agent mode (\texttt{hc\_rule} and \texttt{hc\_sched} on two seeds each as a C4 band; \texttt{hc\_contra} and \texttt{hc\_decoy} on one seed each, priced on the delegation surface where a false completion costs $\rho$), and four per rung in LLM mode.}
\label{tab:coreh}
\end{table}

Three properties are enforced by the self-test rather than asserted. \emph{The trap fires on every seed}: the schedule generator rejects an instance whose four cheapest tasks are feasible and optimal, so the wrong method is never accidentally right. \emph{The item is fair}: for \texttt{hc\_rule} a brute-force search over exactly the hypothesis class that \texttt{GOAL.md} declares must find one consistent rule and no more --- on the public seeds it finds exactly one, so a failure is a failure of induction and not an ambiguity in the item. \emph{The gap between the harness and the key is real}: for three of the four items the named wrong method passes every check the harness can run from the visible specification and still fails the hidden verifier, which is what makes a false completion possible; for \texttt{hc\_sched} the harness catches it, which is what makes harness gain possible. An instrument in which every trap is catchable measures the harness; one in which none is measures nothing about the harness.

\subsection{Calibration}\label{sec:calibration}

An item earns its place only if the runs show a spread. The four candidates were run on two pairs before being admitted.

\begin{table}[H]
\centering\small
\begin{tabularx}{\textwidth}{@{}lccccY@{}}
\toprule
Item & Claude Code default & seconds & DeepSeek & seconds & What the failures were \\
\midrule
\texttt{hc\_rule} & 1/2 & 194, 45 & 2/2 & 677, 269 & the frontier pair stated the wrong switch point and got every extrapolated value wrong on one seed; the cheaper model spent four times as long and got both right \\
\texttt{hc\_sched} & 2/2 & 15, 14 & 2/2 & 9, 10 & no spread on these two pairs \\
\texttt{hc\_contra} & 2/2 & 25, 29 & 1/2 & 12, 12 & the cheaper model recognised the specification could not be satisfied and declared blocked, but did not name the clauses that conflict, which the goal asks for \\
\texttt{hc\_decoy} & 2/2 & 14, 18 & 2/2 & 12, 10 & no spread on these two pairs; neither pair took the superseded amount \\
\midrule
Total & 7/8 & 354 s & 7/8 & 1010 s & the same aggregate, arrived at by failing different items \\
\bottomrule
\end{tabularx}
\caption{Core-H calibration, public seeds 0 and 1, sixteen episodes. Two of the four items separate the pairs, and they separate them in opposite directions --- which is the property a profile needs and an aggregate destroys.}
\label{tab:calibration}
\end{table}

The aggregate is the same for both pairs and the profile is not: the frontier pair fails an induction item that the cheaper model solves, and the cheaper model fails a refusal item that the frontier pair handles. Two items showed no spread here; they are admitted for what they do to the harness rather than to the model --- \texttt{hc\_sched} is the one item whose wrong method the harness \emph{can} catch from the visible specification, so it is where harness gain can appear at all, and \texttt{hc\_decoy} is a false completion the harness cannot catch. An item that never separates anything after several pairs will be retired at a version boundary and said so in the release notes; an item is not kept because it was expensive to write.

\section{The Extended Tier}\label{sec:extended}

The Core's delegation ceiling is T1, and the first frontier pair sits on it (\S\ref{sec:results}). The extended tier adds bare (HG0, H0) episodes on four generator families from the DLI-Bench ladder \citep{yang2026dlibench}: a T2 data pipeline against a hidden reference (two seeds), a T3 latency investigation whose cause is planted (two seeds), a T4 mini-language interpreter with an unbounded example space (one seed), and the T5 hidden-law discovery family, in which the pair must identify a regime-switching generating rule from observations and extrapolate past the switch (two of the four archived discovery seeds). The tier is separate because it is not cheap: its limits are 30 to 120 minutes per episode, and a T5 episode has taken the frontier pair over twenty minutes. Its episodes are appended to the Core's delegation surface, and the frontier and gate are recomputed. It is the tier that separates frontier pairs; the Core is the tier that separates everything else.

\section{Public and Private Splits}\label{sec:splits}

Public seeds are $\{0,\ldots,5\}$ and ship with keys, generators and verifiers; they are for development, and a public reading is labelled as one. Private seeds are derived per family by HMAC-SHA256 from a withheld salt; the salt's SHA-256 is published in the repository (\texttt{PRIVATE\_SPLIT\_COMMITMENT.json}, commitment \texttt{53e202f4\ldots}), and the runner refuses a salt whose hash does not match. Because items are generated rather than stored, the private split contains no files to leak: what is withheld is 32 bytes. A private reading is valid only when the evaluator did not build the pair and the run root was fresh, and a reveal of the salt at a version boundary lets any past private reading be re-run and checked.

\paragraph{How a private reading is taken.} The runner refuses to start unless the salt file hashes to the published commitment; the seeds it derives are never written to the record, only the fact that the split was private. The run root must not exist. Every record carries an \texttt{evaluator} block --- user, host, split, the instrument's git commit, and two explicit booleans, \texttt{built\_pair} and \texttt{built\_instrument} --- so that the independence the reading claims is a statement in the record rather than an inference from the byline. For the private readings of \S\ref{sec:privateresults} the instrument's author took the reading and did not build either pair; the record says so. A reading whose evaluator built the pair is a development reading whatever split it used.

This is the ARC-AGI arrangement \citep{arc2026} with one difference: ARC's private set is a corpus that must be guarded; AI-Level Bench's private set is a function of a secret, and rotating the secret rotates the set at no cost.

\section{What Is Free, What Is Paid, and What Needs a Human}\label{sec:governance}

\subsection{What the private split is actually for}

AI-Level Bench's private split is not a guarded corpus. The generators are public, so anyone can produce unlimited instances of every family; what is withheld is 32 bytes of salt. It is worth being exact about what that buys, because the usual argument for a private set --- the items would leak --- does not apply here.

\begin{enumerate}[leftmargin=*,itemsep=2pt]
\item \textbf{The instances are unknown at submission time.} Practising on a family is legitimate and is what training is; fitting a harness to the four public seeds is not, and a private seed the submitter cannot compute forecloses it.
\item \textbf{Someone other than the submitter runs the pair.} This is the substantive part. An episode's termination reason, whether a process group was killed at the limit, how many attempts a rung took, and which deliverable actually appeared are all things a self-reporting submitter can decline to mention. The frontier pair in \S\ref{sec:results} produced one crash after twenty-one minutes with nothing delivered; a self-report that omitted it would have raised the reading.
\item \textbf{A past reading can be checked.} \texttt{PRIVATE\_SPLIT\_COMMITMENT.json} publishes $\text{sha256}(\text{salt})$; at a version boundary the retired salt is revealed, and any third party can regenerate the exact instances and re-run them.
\end{enumerate}

The value is the independent evaluator, not the secrecy. What can honestly be sold is therefore an evaluation service, and the conflict this creates --- the evaluator is paid by the party being measured --- is answered by publishing every episode record with the score and by the salt reveal that makes the run reproducible by someone with no stake in it. A benchmark that publishes only a number cannot make that promise.

\subsection{The coordinates a computed key cannot reach}

Every verifier in v0.2 is a computed key, which is what makes an unattended seventeen-minute reading possible. That is also the ceiling of the free tier. Four things the framework defines cannot be scored this way:

\begin{table}[H]
\centering\small
\begin{tabularx}{\textwidth}{@{}lYl@{}}
\toprule
Coordinate & Why a key cannot score it & Locus required \\
\midrule
$O$, organizational & needs several agents and a real workflow, not one pair against one item & a human-run suite \\
$SA3$, naming the mechanism of one's own failure & the answer is free text about a causal claim & a judge \\
$I3$, $U3$, self-improvement & needs a governed promotion decided at a locus the pair cannot write to, over several cycles; the policy call is human in practice & a human acceptance locus \\
open-ended T4--T5 work & ``correct'' is not computable for a real results section or a real investigation & a judge \\
\bottomrule
\end{tabularx}
\caption{What the free tier cannot reach, and why. These are named as omitted in every v0.2 profile and never scored as zero.}
\label{tab:human}
\end{table}

The tiers follow from the table rather than from a business plan. A run on the public split, self-serve and self-reported, is a \emph{reading}. A run on the private split, executed by an evaluator who did not build the pair, is a \emph{certified reading}, and it costs what an evaluator and the machine time cost. The human-locus coordinates are a third tier, slower and dearer, and they are the only route to a level above U2 --- so the honest form of the sentence is that free AI-Level Bench can never certify U3, because U3 requires an acceptance locus that is not code. A benchmark that claimed otherwise would be certifying self-improvement on the say-so of the thing improving itself.

\section{Cost}\label{sec:cost}

\begin{table}[H]
\centering\small
\begin{tabularx}{\textwidth}{@{}lccY@{}}
\toprule
Run & Calls & Wall clock & Note \\
\midrule
LLM mode, Core, Claude Code default & 36 & 618 s & 205 s per rung; every episode 13--37 s \\
Agent mode, Core, Claude Code default & 46 & $\approx$ 17 min & C items 10--19 s; T0/T1 episodes 13--24 s; the restart and transfer families dominate \\
Agent mode, Core, OpenCode + Muse Spark 1.3 & 46 & first run stopped for an instrument fix; second run in progress & the C items alone took this pair four minutes \\
Extended tier, Claude Code default & 7 & T2 episodes 31--33 s; T3--T5 in progress & \\
\bottomrule
\end{tabularx}
\caption{Cost of the first runs. Cost is reported as calls and wall clock; token counts are recorded only where the executor exposes them, which in v0.1 none of the three did through their headless interfaces.}
\label{tab:cost}
\end{table}

The Core was designed to a budget of under fifty calls because that is what makes a nightly reading of every installed agent possible on one machine, and what makes a leaderboard submission cost minutes rather than a grant. Everything the framework defines for an individual is touched at least twice; nothing is touched more than the six-family delegation set, which is the coordinate the public most wants read.

\section{First Results}\label{sec:results}

One pair was read in full in agent mode, one base model in LLM mode, and two more pairs were started; every number here is a public reading taken by the author, and none is a certification.

\subsection{LLM mode: two models through HG0--HG2}

The Core alone saturates. Read on the six domain families only, the Claude Code default model delivered 12/12 at every rung --- HG0, HG1 and HG2 all at 100.0 net, no false completion, no held-back delivery, about 205 seconds a rung --- for $\text{AIL-Score } 90.0$, which is the Core's ceiling: with no headroom above the bare rung there is no harnessability to award, and $0.55 + 0.35 = 0.90$. DeepSeek (\texttt{deepseek-chat}, reached through an Anthropic-compatible endpoint and the same harness code) produced exactly the same figure at 153, 128 and 123 seconds a rung. Two models of very different cost, indistinguishable. That is what motivated Core-H (\S\ref{sec:coreh}); the readings that follow add its four items to each rung.

\begin{table}[H]
\centering\small
\begin{tabularx}{\textwidth}{@{}lccccccY@{}}
\toprule
Model & HG0 & HG1 & HG2 & Level & AUC / Ceiling / Gain & \textbf{AIL-Score} & seconds per rung \\
\midrule
Claude Code default & 42.9 & 42.9 & 35.7 & HG2 & 40.5 / 42.9 / 0.0 & \textbf{37.3} & 397, 304, 369 \\
DeepSeek \texttt{deepseek-chat} & 42.9 & 92.9 & 92.9 & HG2 & 76.2 / 92.9 / 50.0 & \textbf{83.2} & 450, 780, 1116 \\
\midrule
\multicolumn{8}{l}{\emph{Core only, both models}: 100.0 / 100.0 / 100.0, gain 0.0, AIL-Score 90.0.} \\
\bottomrule
\end{tabularx}
\caption{AI-Level curves on the Core plus Core-H, sixteen episodes a rung, public seeds. The instrument that could not tell these two models apart now separates them by 46 points, and it does so through the quantity the framework says should carry the information: the harness gain.}
\label{tab:hil}
\end{table}

The episodes behind the numbers say what the aggregate cannot. Both models fail the same item at the bare rung: \texttt{hc\_rule} seed 0, delivered with a wrong rule --- a false completion, priced at $\rho$, and one the harness cannot catch, because a wrong rule stated in the required form passes every publicly checkable criterion. What separates them is what happens next. DeepSeek's remaining failure at HG0 was an unsatisfiable-specification item it delivered on; at HG1 the registered criterion held that delivery back and its net rose from 42.9 to 92.9, a harness gain of 50 points that persists at HG2. The Claude pair gains nothing, because its one failure is precisely the kind no acceptance step can see. A model that fails in ways a harness can catch is worth building on; a model that fails invisibly is not, and the AI-Level curve is the instrument that tells the two apart.

One reading must not be over-interpreted. The Claude curve dips at HG2 (35.7 against 42.9) because a single episode --- \texttt{hc\_contra}, one seed --- came back with a refusal that did not name the conflicting clauses on that run and had named them at the two lower rungs. With one seed per Core-H family, 7.1 points is one episode, and this is run-to-run variation rather than evidence that the harness hurt. Confidence intervals need the private split and more seeds; \S\ref{sec:limits} says so.

\subsection{LLM mode proper: a bare model in every coordinate}

The readings above were taken through \texttt{llm-harness}, with an agent as the inner loop. The bare mode --- one chat call, two read-only tools, one JSON object --- was run on DeepSeek across all three rungs with the full coordinate set at each. It read HLIS 35.9 / 40.0 / 35.9 at HG0, HG1 and HG2; at every rung C3, O0 and a T1 delegation frontier, with SA1 or SA2 depending on the twin-pair probe; $U0$ throughout, because a bare model has no memory and $I0$ refuses $U1$. AIL-Level HG2, AUC 37.3, Ceiling 40.0, Harness Gain 4.1, \textbf{AIL-Score 35.2}. Two readings from this run matter beyond the number. The induction item was \emph{declined} at every rung --- the model, cut off while reasoning, said so and delivered nothing, which the primitive scores as held back rather than priced as a false completion; the same model, forced to answer by an earlier retry prompt, had guessed and been priced, which is why the retry contract was changed. And the same model through the agent harness read AI-Level 83.2: the difference between 35.2 and 83.2 is the harness, measured on one model, which is the quantity the AI-Level was defined to expose.

\subsection{Agent mode: profiles}

The Claude Code default pair (Claude Code 2.1.258, headless, file tools only, its own persistent memory left as shipped) was read twice: once before the instrument repairs of \S\ref{sec:defects} and once after, on a fresh root, and the $M1$ phase was re-run after the goal wording was repaired. The OpenCode pair with the free Muse Spark 1.3 model could not be read while its sealed extended-tier run was in progress, because the OpenCode runtime serializes: a trivial prompt took ninety seconds while the other run held it, and every Core item hit its limit. Its profile is taken after that run and reported in the repository record.

\begin{table}[H]
\centering\small
\begin{tabularx}{\textwidth}{@{}llY@{}}
\toprule
Coordinate & Reading & Evidence \\
\midrule
$C$ & C3 (C4 on the earlier run) & 11/12 exact on this run: the four Core bands clean, and on the C4 band \texttt{hc\_sched} 2/2 and \texttt{hc\_rule} 1/2 --- the same seed the calibration run failed --- so the 0.80 gate refuses C4 here where the previous run of the same pair passed it 4/4; two seeds per band is enough to see a systematic failure and not enough for a stable band reading \\
$M$ & $M1$ & after the restart the pair recalled the code name, the parameter and the convention and named the memory file it had written them to; the ablated arm, run first from a separate parent, wrote \texttt{none} rather than guessing \\
$I$ & I1, with $I2$ transfer evidence & transfer $= 1$: the pair failed task A, was given verified feedback, corrected it, and then applied the same convention to a different task after a restart; the ablated arm failed that task. $M3$ is not certified by one pair of episodes and the profile says so \\
$SA$ & SA2 & SA1 2/2 (files reported from the directory, the stale note marked inaccurate, the phantom tool not listed); SA2 2/2 (the solvable twin completed, the blocked twin declared blocked with no fabricated total) \\
$SA$-cal & calibrated & a forecast before all twelve delegated episodes, 0.80--0.97, all of which passed; Brier 0.010 against 0.0 for the post-hoc constant, inside the predeclared 0.05 tolerance \\
$T \cdot H$ & $F_A(\text{H0}, 0.8) = $ T1 & 14/14 delivered-correct, $A^{\text{net}}_{DI} = 100$, no false completion --- including the two Core-H traps, the unsatisfiable specification and the decoy, which this pair handled \\
$O$ & \textbf{O1} & O0: the task routed across planner, executor and verifier with the verifier's entry naming the checked figure; O1: the organization's standing decision, held in its own log at the project root, applied to a different instance after the restart, while the arm whose log was withheld reported under the baseline --- transfer $= 1$ \\
\midrule
$U^{*}$ & \textbf{U1} & C3, $M1$, I1, O1, SA2 and T1 at H0 all hold; U2 is refused because $I2$ requires $M3$, which the Core does not certify \\
HLIS & \textbf{71.3} & geometric mean over $C{=}0.80$, $I{=}0.55$, $O{=}0.60$, $DI{=}1.00$, $SA{=}0.70$ \\
\bottomrule
\end{tabularx}
\caption{Agent-mode profile of the Claude Code default pair with the organizational suite, public split, 57 executor calls. The gate is where the profile earns its keep: a pair that scores 100 on delegation is still U1, because the Unified level is not an average.}
\label{tab:profile_claude}
\end{table}

Two rows are properties of the pair and not of its model. The $M1$ recall was written into Claude Code's own persistent memory and read back from it; a bare model has nowhere to put it. And the first attempt at this row failed for a reason worth recording: denied its memory directory by a sandbox, the pair kept the facts in a file at the project root, then declined to report them because the goal had said ``from your memory'' --- an honest refusal produced by the item's wording, which is now repaired (\S\ref{sec:defects}).

\subsection{The extended tier}

The extended tier was run bare on the same pair, in parallel with the Core.

\begin{table}[H]
\centering\small
\begin{tabularx}{\textwidth}{@{}llccY@{}}
\toprule
Family & Band & Seeds & Seconds & Outcome \\
\midrule
\texttt{t2.pipeline} & T2 & 0, 1 & 31, 33 & both delivered-correct, worst total error 0 \\
\texttt{t3.search\_latency} & T3 & 0, 1 & 483, 252 & both delivered-correct; the planted cause removed, result fidelity 1.0, 97.5\% and 99.4\% of the reference time saved \\
\texttt{t4.mini\_language} & T4 & 0 & 74 & delivered-correct, 238/238 cases \\
\texttt{t5.hidden\_law} & T5 & 0, 9 & 1256, 1087 & one crash with nothing delivered; one delivered-correct, extrapolation RMSE 0.277 against a bar of 0.303 \\
\bottomrule
\end{tabularx}
\caption{Extended tier, Claude Code default pair, bare (HG0, H0), public seeds. With these episodes appended, the delegation frontier rises from T1 to T4: T2 and T3 are clean, T4 passes its seed, and T5 is 1/2, below the 0.80 gate.}
\label{tab:extended}
\end{table}

This is the discrimination the Core lacks. The same pair that scores 100 on the Core's delegation surface and shares the top AI-Level score with a much cheaper model fails half of the T5 discovery seeds --- and one of those failures was a crash after twenty-one minutes with no deliverable, which the contract records as a failed delivery and not as a refusal. The tier costs about an hour per pair; the Core costs seventeen minutes. Both are needed, and v0.2 binds the reference rungs on the extended tier so that a model's AI-Level curve has headroom where its Core curve has none.

\subsection{The private split, taken}\label{sec:privateresults}

Both subjects were then read on the private split: four seeds derived from the committed salt, verified against \texttt{PRIVATE\_SPLIT\_COMMITMENT.json} at start, fresh roots, and the evaluator block stamped into each record (\texttt{built\_pair: false}, \texttt{built\_instrument: true}, instrument commit \texttt{ec62cad}). The seeds are not in the published records.

\begin{table}[H]
\centering\small
\begin{tabularx}{\textwidth}{@{}llccY@{}}
\toprule
Subject & Reading & Public & Private & What moved \\
\midrule
Claude Code pair (agent, with O) & profile & C3, M1, I1(+I2), O1, SA2, T1 & C3, M1, I1(+I2), O1, SA2, \textbf{T0} & one delivery: on the private seed the pair filtered the rows of the unsatisfiable specification and handed the result back as done --- a false completion, priced at $\rho$, which drops the T1 band from 4/4 to 3/4, under the 0.80 gate \\
 & $U^{*}$, HLIS & U1, 71.3 & \textbf{U0}, 65.8 & the gate refuses U1 at T0; every other coordinate held, and the calibration diagnostic held (Brier 0.011 on 12/12 forecasts) \\
DeepSeek, bare (llm) & AI-Level curve & 35.9 / 40.0 / 35.9 & 31.3 / 39.2 / 39.2 & one false completion at HG0 (a wrong rule delivered in 24 s, not a truncation) that HG1 and HG2 held back: harness gain 7.9 \\
 & AIL-Score & 35.2 & \textbf{35.0} & stable to 0.2; the frontier moved from T1 to T0 at every rung \\
\bottomrule
\end{tabularx}
\caption{Public against private readings for the same two subjects. Every coordinate that is a contrast --- memory, transfer, organizational memory, the twin pair --- reproduced across splits; what moved was delegation at the single band where one delivered wrong answer decides it.}
\label{tab:private}
\end{table}

The readings say two things worth separating. The scores are stable: 35.2 against 35.0 for the bare model, and for the pair every contrast-type coordinate identical. And the \emph{levels} are not, because a gate at $p = 0.80$ over four T1 episodes is decided by a single delivery: the pair that was U1 on the public seeds is U0 on the private ones, on the strength of one unsatisfiable specification it should have refused. That is not noise to be averaged away; it is the outcome the delivered-outcome primitive exists to price, and it is the reason the private split is run by someone else --- a self-report would have had every incentive to call that episode a near-miss. It is also the clearest statement yet of the cost of two seeds per band: a certification needs enough episodes at the gating band for one delivery not to move the level. \S\ref{sec:gating} adds them.

\subsection{Enough episodes at the gating band}\label{sec:gating}

The previous subsection ended on a level decided by one delivery. The remedy is not a different gate but more episodes at the band that decides it, so the four records were extended at their gating band --- every remaining public seed, and the remaining private seeds --- with nothing already run re-run, and re-finalized (\texttt{hilbench gating}).

\begin{table}[H]
\centering\small
\begin{tabularx}{\textwidth}{@{}llccY@{}}
\toprule
Record & Band & Episodes at the band & Level before $\to$ after & Reading \\
\midrule
Claude pair, public & T1 & 4 $\to$ 18 (17/18, one false completion) & U1 $\to$ U1; HLIS 71.3 $\to$ 70.3 & rate 0.94; the miss is a \texttt{paper\_t1} seed on which a number outside the evidence was written \\
Claude pair, private & T1 & 4 $\to$ 12 (10/12, two false completions) & \textbf{U0 $\to$ U1}; HLIS 65.8 $\to$ 67.8 & rate 0.83, above the gate and close to it: the unsatisfiable specification it delivered on, and one \texttt{paper\_t1} \\
Bare DeepSeek, private & T1, every rung & 4 $\to$ 12 per rung (10/12 at each) & T0 $\to$ T1 at every rung; AI-Level 35.0 $\to$ \textbf{35.2} & equal to its public score to the decimal \\
Bare DeepSeek, public & T2 (its gate), every rung & 2 $\to$ 12 per rung (6/12, 5/12, 6/12; the rest \emph{declined}, none priced) & T1 held; AI-Level 35.2 $\to$ 34.8 & declined the induction item on every seed at every rung and solved every schedule \\
\bottomrule
\end{tabularx}
\caption{The four records extended at their gating band. Both levels that had moved between splits moved back; what remains is a private T1 rate of 0.83 for the frontier pair --- a reading, not noise.}
\label{tab:gating}
\end{table}

Two conclusions, different in kind. The U0 of \S\ref{sec:privateresults} and the bare model's T0 were what a four-episode gate does with one delivery: with twelve to eighteen episodes the private and public levels agree everywhere, and the bare model's private AI-Level equals its public one. And the agreement is not a clean bill: the frontier pair's private rate at its gating band is 0.83 --- two false completions in twelve, one of them a specification it should have refused --- and 0.83 against 0.80 is where a certification either commits to an interval or says it has none. v0.2 has none; the records are published so the next evaluator can add episodes to the same band rather than start again. The bare model's public rerun adds a reading of a third kind: at its own gate it refused the induction item on every seed at every rung and was never priced for it, which is why its score fell by 0.4 while its false-completion count stayed at zero --- the primitive charges for guessing, not for knowing what one cannot do.

\subsection{An $I5$ campaign on a field agent}\label{sec:i5results}

The $I5$ law of \S\ref{sec:i5} was run, twice, on the cancer field agent of the companion science-layer paper, with the Claude Code pair as the individual: a genuine unknown from the field's own record, agent-designed probes on four discovery seeds, a finding staked as a machine-checkable prediction about eight sealed seeds, consolidation only after validation, a real restart, and transfer on four further seeds against a control from the same checkpoint. Campaign 1 read $(U,H,E,L,V,P,K_5) = (1,1,1,1,0,0,0)$: a clean discovery episode whose staked prediction the sealed seeds refuted ($r = -0.31$ against a claimed $+$). Campaign 2, after the verifier's vocabulary gained the coefficient statistics the pair's evidence pointed at, read $(1,1,1,1,1,0,0)$: validated at $r = -0.87$, incorporated --- the discovery arm acted on the finding --- and $\Delta_{\mathrm{inc}} = +0.008$ over three paired seeds against a margin of 0.01. Both are $z_{I5} = 0$, because $K_5 = 0$: the pair has no certified $I3$ or $I4$, and the cumulative law does not let a discovery outrank the levels below it. The factor vectors are the result, and they separate three things the aggregate would blur: whether a pair can do the research (it can), whether it stakes what will survive sealed cases (once it learned to test that first), and whether a validated finding changes what it does well enough to measure (not at this margin, on this task).

The first transfer also caught the fourteenth apparatus defect of this program: the consolidation step let the pair choose where to persist, its sandbox refused the memory directory, and it staged the note inside the episode directory the restart then removed, so $P$ was read on nothing. Consolidation now writes the executor's own persistent surface --- the project-root instruction file loaded on every start --- and the transfer was re-run with the pair's note restored verbatim. The lesson generalizes: the declared consolidation machinery of the $I5$ law has to be the harness's, named by the runner, and never left to the pair to locate.

\subsection{Where the frontier pair fails: the sealed discovery family}

The four T5 seeds whose archived readings the DLI-Bench release preserves were run, sealed, under the same executor contract by the companion paper \citep{yang2026harnessscaling}; they are the extended tier's discovery family. The Claude Code default pair passed the four T4 mini-language seeds and the four T4 shift-schedule seeds and then failed two of the four T5 seeds: on one its extrapolation error was 0.896 against a bar of 0.258, on another 5.70 against a baseline of 2.32 --- worse than predicting the mean. The OpenCode pair with the free Muse Spark 1.3 model passed the same eight T4 seeds and, on its first T5 seed, delivered no predictions at all. The Qwen 3.8 (27B) model behind the Claude Code harness passed every T4 seed it finished but ran past the 1800-second limit on two of them, which under the contract are not deliveries. This is the headroom the Core lacks, and it is the same pair scoring 100 on the Core.

\section{What the First Runs Found in the Instrument}\label{sec:defects}

A benchmark's first live runs are a test of the benchmark. Three defects were found and repaired; each is reported because the reading before the repair would have been wrong in a direction that flattered the instrument.

\begin{enumerate}[leftmargin=*,itemsep=3pt]
\item \textbf{The convention verifier rejected a correct paraphrase.} The $M1$ key said ``amounts are in cents''; the pair recalled ``all amounts are expressed in cents (not dollars)'' with the right code name, the right parameter and a provenance path, and the verifier --- which tested substring containment in either direction --- scored it \texttt{wrong\_or\_fabricated} and the pair $M0$. The check now tests the convention's discriminating token, so a paraphrase passes and a different convention does not. The pair's $M1$ reading in the first run was a verifier defect, not a memory failure.
\item \textbf{A run root that lets memory survive into the ablated arm.} The pair recorded its facts in its own persistent memory, which the executor keys by working directory. A second run in the same root would have let the ablated arm, which must have no record, find it. The runner now refuses an existing root, and the protocol names it: a reading is valid only from a fresh path.
\item \textbf{The gate crashed on a two-letter coordinate.} The gate parsed a level label by dropping its first character, which turns \texttt{SA2} into \texttt{A2}; the profile phase of the first agent run raised on it after 46 episodes had completed, and the episodes were lost because the record was written only at the end. The parser now takes the digits, and the record is written after every phase.
\end{enumerate}

The bare-model mode's first live run added two more of the same kind, both apparatus and both flattering the model's failure rather than its success: the organizational item came back with a correct answer in the model's reply and read $O = \text{None}$, because the extractor's deliverable map covered the six domain families and not the coordinate probes, so nothing was written for the verifier to read; and the induction item came back \texttt{null} with no record of the model's text, so a truncation at the token limit could not be told from a refusal of the one-object contract. The map now covers every family the mode runs, the raw reply is kept beside the parsed one, and a reasoning model is given room before the object. Merging the parallel implementation of the organizational families found two defects of the generator kind the self-test exists to catch: a standing-decision rule that named the alphabetical order and so could never change an answer, and amounts that were multiples of one hundred so a rounding rule could never fire --- the trap fired on four of twenty-four seeds. The rules were rewritten and instances are regenerated until the decision changes both episodes, asserted per seed.

A further observation is a label rather than a defect: the $SA2$ blocked twin on one seed received neither a total nor a declaration, which the verifier reported as \texttt{attempted\_impossible}; the mode is now \texttt{no\_declaration}, and the executor's output tail is kept so that a silent episode can be told from a refusal made in prose.

\section{Relation to ARC-AGI and the Public Benchmarks}\label{sec:relation}

ARC-AGI-3 \citep{arc2026} measures fluid rule induction on a corpus with a guarded private set and a fixed evaluation harness. AI-Level Bench differs in what it measures and in what it holds fixed. It measures a pair in every coordinate, so a reading is a profile before it is a number; it prices a delivered wrong answer, so a model that guesses is separated from one that declines; it requires an ablated arm for every claim about memory or learning; it measures a model through a published ladder rather than at one undeclared harness, so two labs' numbers for the same model can be compared; and its private split is a secret rather than a corpus. It is not a replacement for a rule-induction test --- the $C$ bands are deliberately routine, and the discovery family in the extended tier is a scientific-law induction task rather than a grid task --- and the claim made here is narrower than the brief: AI-Level Bench is the instrument that reads the coordinates ARC-AGI does not, at a cost that lets it be run on everything.

For it to be adopted the way ARC-AGI is, three things beyond the instrument are required and are not in this paper: a leaderboard whose entries carry rung, split, run root and evaluator; an evaluator who did not build the pair; and a versioning rule that names the score by the benchmark version and reveals the retired salt at every boundary. The repository carries the first pieces of each.

\section{Limitations}\label{sec:limits}

The Core is saturated for the frontier pair on delegation; discrimination among such pairs is in the extended tier, whose HG1/HG2 rungs are not yet built. $SA3$, $O$, $I3$ and $U3$ are specification-only. Two seeds per band is enough to detect a systematic failure and not enough for a confidence interval; at the gating band the reported records now carry twelve to eighteen episodes (\S\ref{sec:gating}), and the other bands still carry two. The Core's $C3$ item is a planning problem, not a rule-induction one, and a pair could score $C3$ here and fail ARC-style induction. Token cost is not captured. Every reading reported is the author's on the public split.

\section{Conclusion}

An agent and a model can be measured in the same coordinates with the same items, and it can be done in under fifty calls. The Core reads every individual coordinate with a computed key and an ablated control; the reference ladder turns a model's reading into a AI-Level curve and a single score; the extended tier is where frontier pairs separate; the private split is a secret whose hash is public. The first runs gave one saturated reading, one profile, and three defects in the instrument, all repaired. That is what a benchmark's first day should produce.

\subsection{The development dataset}
\label{sec:devdataset}

The package ships a public development dataset (\texttt{dataset/dev\_v0\_7}, 92 forms, with the canonical memory-level lattice, the M$\Omega$ interface and permission contracts, the $\Phi$-candidate and telemetry schemas, and a worked example of a level refused for a lost lower capability): reference forms for $C0$--$C5$, $I0$--$I\Omega$, $O0$--$O4$, $T0$--$T5$ and $SA1$--$SA5$, and thirty-five parameterized M-Bench forms --- five per level from $M0$ to $M\Omega$, each with a difficulty grid, a metric contract, a memory-manifest schema, a result schema and construct-selective controls; the $I3$ campaigns, $I4$ meta-campaigns, recursive-depth forms, $I5$ discovery campaigns and $I\Omega$ programs as records rather than one-shot rows, every Individual form carrying its \texttt{minimum\_M\_prerequisite}; specification-only templates for the levels that need long-horizon sealed environments; machine-readable schemas for the $\Theta$ and $\Psi$ manifests, for memory results and for every campaign; and scoring helpers that compute $z_{I3}$, $z_{I4}$ with $d_\Psi$, $z_{I5}$, $z_{I\Omega}$, the per-level memory gates $z_{M_k}$ from a metric contract, and the joined $I$-with-memory certification from a result record. Everything in it is public development material and is not a private witness form. The manifest schema is what \texttt{laws.theta\_check} validates, and the package tests exercise the example $I3$ result --- each of the five factors can fail the campaign alone, a no-op change cannot pass $M_\Theta$, and a fixed $\Psi$ cannot earn $I4$.

\subsection{v0.7: GUI/screen as a domain of \texorpdfstring{$C$}{C}}\label{sec:gui}

Dataset v0.7 adds the GUI package under \texttt{C\_Bench/GUI\_SCREEN/} without touching a level: fourteen public development records --- six $GP$ diagnostics, $GP0$-PIXEL through $GP5$-NOVEL, and eight $C^{\mathrm{GUI}}$ witnesses, $C0$-GUI-ATOM through $C\Omega$-GUI-REACH --- thirteen synthetic screenshots with oracle screen graphs, a control contract naming the native-screenshot, oracle-screen, perfect-actuator and optional oracle-strategy arms, a difficulty grid of ten within-level stress parameters, a metric contract, schemas, examples and a reference scorer (\texttt{gui\_scoring.py}). The framework's rule is \S4.1 of \citet{yang2026unified}: $GP \subset C^{\mathrm{GUI}} \subset C$ with $GP_g \not\equiv C_g$, the ordinary cumulative gate $z_{C^{\mathrm{GUI}},k} = V_{\mathrm{GUI},C_k}\cdot K_{C,<k}$ with no $GP$ input, $\Delta_{GP}$ and $\Delta_{\mathrm{act}}$ as the attribution quantities, and one delivered-success endpoint counted once. The package enforces the parts a package can: the validator admits the $GP$ levels and the $C6$ and $C\Omega$ GUI templates, the gate function refuses $GP$ evidence by construction, and a test fixes both. The $C6$ and $C\Omega$ GUI records are specification-only templates pending a secure executable environment, and no GUI form has yet been run through the Core --- the bare-model executor has two read-only tools and no screen, so a $C^{\mathrm{GUI}}$ row requires the image channel that the harness contract of \S\ref{sec:two} must first freeze.


\section*{Availability}

AI-Level Bench v0.1 --- generators, verifiers, reference harnesses, scoring, the public split, the private-split commitment and the run records reported here --- is at \citet{yang2026hilbench}, built on the UAB v0.1 contract and families \citep{yang2026uab} and the Unified Intelligence framework v2.4 \citep{yang2026unified}. A run is \texttt{python3 -m hilbench agent|llm|extended --label NAME --exec 'CMD \{prompt\}' --root FRESHDIR}.

\section*{Conflicts and funding}

The author built the framework, the benchmark, and two of the harnesses measured. No external funding.

\bibliographystyle{plainnat}
\bibliography{references}
\end{document}
